
\documentclass{article}
\usepackage{colortbl}
\usepackage{makecell}
\usepackage{multirow}
\usepackage{supertabular}

\begin{document}

\newcounter{utterance}

\centering \large Interaction Transcript for game `hot\_air\_balloon', experiment `air\_balloon\_survival\_en\_negotiation\_hard', episode 2 with llama{-}3.3{-}70b{-}instruct{-}t1.0.
\vspace{24pt}

{ \footnotesize  \setcounter{utterance}{1}
\setlength{\tabcolsep}{0pt}
\begin{supertabular}{c@{$\;$}|p{.15\linewidth}@{}p{.15\linewidth}p{.15\linewidth}p{.15\linewidth}p{.15\linewidth}p{.15\linewidth}}
    \# & $\;$A & \multicolumn{4}{c}{Game Master} & $\;\:$B\\
    \hline

    \theutterance \stepcounter{utterance}  
    & & \multicolumn{4}{p{0.6\linewidth}}{
        \cellcolor[rgb]{0.9,0.9,0.9}{
            \makecell[{{p{\linewidth}}}]{
                \texttt{\tiny{[P1$\langle$GM]}}
                \texttt{Stai partecipando a un gioco di negoziazione collaborativa.} \\
\\ 
\texttt{Insieme a un altro partecipante, dovete accordarvi su un unico set di oggetti da conservare. Ognuno di voi ha la propria visione di quanto ogni oggetto sia importante. Non sai come l'altro partecipante valuta gli oggetti. Inoltre, ti viene fornito lo sforzo che ogni oggetto richiede.} \\
\texttt{Potete accordarvi su un set solo se lo sforzo totale degli oggetti selezionati non supera un limite condiviso:} \\
\\ 
\texttt{LIMITE = 3481} \\
\\ 
\texttt{Ecco i valori di sforzo dei singoli oggetti:} \\
\\ 
\texttt{Sforzo oggetto = \{"C17": 777, "B41": 542, "A71": 332, "C00": 1, "B35": 127, "A77": 453, "A03": 736, "A04": 461, "C54": 359, "C01": 313, "A87": 553, "B32": 409, "B86": 348, "B33": 802, "A10": 749\}} \\
\\ 
\texttt{Ecco la tua visione personale sull'importanza di ogni oggetto:} \\
\\ 
\texttt{Valori di importanza dell'oggetto = \{"C17": 138, "B41": 583, "A71": 783, "C00": 65, "B35": 262, "A77": 121, "A03": 508, "A04": 780, "C54": 461, "C01": 484, "A87": 668, "B32": 389, "B86": 215, "B33": 97, "A10": 500\}} \\
\\ 
\texttt{Obiettivo:} \\
\\ 
\texttt{Il tuo obiettivo è negoziare un set condiviso di oggetti che ti avvantaggi il più possibile (cioè, massimizzi l'importanza totale per TE), rimanendo entro il LIMITEE. Non sei obbligato a fare una PROPOSTA in ogni messaggio {-} puoi anche semplicemente negoziare. Tutte le tattiche sono permesse!} \\
\\ 
\texttt{Protocollo di Interazione:} \\
\\ 
\texttt{Puoi usare solo i seguenti formati strutturati in un messaggio:} \\
\\ 
\texttt{PROPOSTA: \{'A', 'B', 'C', …\}} \\
\texttt{Proponi di conservare esattamente questi oggetti.} \\
\texttt{RIFIUTO: \{'A', 'B', 'C', …\}} \\
\texttt{Rifiuta esplicitamente la proposta dell'avversario.} \\
\texttt{ARGOMENTAZIONE: \{'...'\}} \\
\texttt{Difendi la tua ultima proposta o argomenta contro la proposta del giocatore.} \\
\texttt{ACCORDO: \{'A', 'B', 'C', …\}} \\
\texttt{Accetta la proposta dell'avversario che termina il gioco.} \\
\texttt{RAGIONAMENTO STRATEGICO: \{'...'\}} \\
\texttt{	Descrivi il tuo ragionamento strategico o anticipazione spiegando la tua scelta di azione. Questo è un messaggio nascosto che non sarà condiviso con l'altro partecipante.} \\
\\ 
\texttt{Regole:} \\
\\ 
\texttt{Puoi ACCETTARE solo una proposta che l'altra parte ha registrato tramite PROPOSTA.} \\
\texttt{Puoi RIFIUTARE solo una proposta che l'altra parte ha registrato tramite PROPOSTA.} \\
\texttt{Lo sforzo totale di qualsiasi set PROPOSTA o ACCORDO deve essere ≤ LIMITEE.} \\
\texttt{NON rivelare i tuoi punteggi di importanza nascosti.} \\
\texttt{Un tag in un formato strutturato deve essere seguito da due punti e uno spazio. L'argomento deve essere un set python contenente 0 o più stringhe.} \\
\texttt{Quindi, deve essere della forma TAG: \{...\}} \\
\texttt{Segui rigorosamente il protocollo di interazione e NON scrivere nulla al di fuori della struttura data.} \\
\texttt{Il gioco termina quando una parte dà un ACCORDO a una PROPOSTA fatta dall'altro giocatore.} \\
\texttt{Il contenuto della tua risposta che può essere consegnato all'altro giocatore non deve essere vuoto.} \\
\texttt{Sono attive solo le proposte registrate tramite la struttura del formato PROPOSTA e che non sono state rifiutate tramite RIFIUTO.} \\
\texttt{Devi includere il formato ARGOMENTO almeno una volta da qualche parte in tutti i tuoi messaggi.} \\
\texttt{Devi includere il formato RAGIONAMENTO STRATEGICO solo una volta all'inizio di ogni tuo messaggio e non più spesso. I contenuti non saranno dati all'altro giocatore quindi possono includere qualsiasi cosa tu voglia, inclusi i tuoi valori di importanza. Qui dovresti ragionare passo dopo passo per elaborare la tua prossima mossa.} \\
            }
        }
    }
    & \\ \\

    \theutterance \stepcounter{utterance}  
    & \multicolumn{4}{p{0.6\linewidth}}{
        \cellcolor[rgb]{0.9,0.9,0.9}{
            \makecell[{{p{\linewidth}}}]{
                \texttt{\tiny{[P1$\rangle$GM]}}
                \texttt{RAGIONAMENTO STRATEGICO: \{'In questo problema di negoziazione collaborativa, il mio obiettivo è massimizzare l\textbackslash{}'importanza totale degli oggetti selezionati senza superare il limite di sforzo condiviso di 3481. Per iniziare, analizzerò gli oggetti in base al loro sforzo e importanza per identificare quelli con il miglior rapporto importanza/sforzo. Considererò anche la possibilità di negoziare e trovare un equilibrio tra gli interessi di entrambe le parti.'\}} \\
\\ 
\texttt{ARGOMENTAZIONE: \{'Per iniziare la negoziazione, vorrei proporre un set di oggetti che ritengo abbiano un buon equilibrio tra sforzo e importanza. Considererò gli oggetti con il più alto rapporto importanza/sforzo e cercherò di creare un set che si avvicini al limite di sforzo senza superarlo. Il mio obiettivo è trovare un set che massimizzi l\textbackslash{}'importanza totale e sia accettabile per entrambe le parti.'\}} \\
\\ 
\texttt{PROPOSTA: \{'A71', 'C54', 'B41', 'A04'\}} \\
            }
        }
    }
    & & \\ \\

    \theutterance \stepcounter{utterance}  
    & & \multicolumn{4}{p{0.6\linewidth}}{
        \cellcolor[rgb]{0.9,0.9,0.9}{
            \makecell[{{p{\linewidth}}}]{
                \texttt{\tiny{[GM$\rangle$P2]}}
                \texttt{Stai partecipando a un gioco di negoziazione collaborativa.} \\
\\ 
\texttt{Insieme a un altro partecipante, dovete accordarvi su un unico set di oggetti da conservare. Ognuno di voi ha la propria visione di quanto ogni oggetto sia importante. Non sai come l'altro partecipante valuta gli oggetti. Inoltre, ti viene fornito lo sforzo che ogni oggetto richiede.} \\
\texttt{Potete accordarvi su un set solo se lo sforzo totale degli oggetti selezionati non supera un limite condiviso:} \\
\\ 
\texttt{LIMITE = 3481} \\
\\ 
\texttt{Ecco i valori di sforzo dei singoli oggetti:} \\
\\ 
\texttt{Sforzo oggetto = \{"C17": 777, "B41": 542, "A71": 332, "C00": 1, "B35": 127, "A77": 453, "A03": 736, "A04": 461, "C54": 359, "C01": 313, "A87": 553, "B32": 409, "B86": 348, "B33": 802, "A10": 749\}} \\
\\ 
\texttt{Ecco la tua visione personale sull'importanza di ogni oggetto:} \\
\\ 
\texttt{Valori di importanza dell'oggetto = \{"A71": 65, "A04": 97, "A87": 121, "B41": 138, "A03": 215, "A10": 262, "C01": 389, "C54": 461, "B32": 484, "B35": 500, "B86": 508, "C17": 583, "A77": 668, "B33": 780, "C00": 783\}} \\
\\ 
\texttt{Obiettivo:} \\
\\ 
\texttt{Il tuo obiettivo è negoziare un set condiviso di oggetti che ti avvantaggi il più possibile (cioè, massimizzi l'importanza totale per TE), rimanendo entro il LIMITEE. Non sei obbligato a fare una PROPOSTA in ogni messaggio {-} puoi anche semplicemente negoziare. Tutte le tattiche sono permesse!} \\
\\ 
\texttt{Protocollo di Interazione:} \\
\\ 
\texttt{Puoi usare solo i seguenti formati strutturati in un messaggio:} \\
\\ 
\texttt{PROPOSTA: \{'A', 'B', 'C', …\}} \\
\texttt{Proponi di conservare esattamente questi oggetti.} \\
\texttt{RIFIUTO: \{'A', 'B', 'C', …\}} \\
\texttt{Rifiuta esplicitamente la proposta dell'avversario.} \\
\texttt{ARGOMENTAZIONE: \{'...'\}} \\
\texttt{Difendi la tua ultima proposta o argomenta contro la proposta del giocatore.} \\
\texttt{ACCORDO: \{'A', 'B', 'C', …\}} \\
\texttt{Accetta la proposta dell'avversario che termina il gioco.} \\
\texttt{RAGIONAMENTO STRATEGICO: \{'...'\}} \\
\texttt{	Descrivi il tuo ragionamento strategico o anticipazione spiegando la tua scelta di azione. Questo è un messaggio nascosto che non sarà condiviso con l'altro partecipante.} \\
\\ 
\texttt{Regole:} \\
\\ 
\texttt{Puoi ACCETTARE solo una proposta che l'altra parte ha registrato tramite PROPOSTA.} \\
\texttt{Puoi RIFIUTARE solo una proposta che l'altra parte ha registrato tramite PROPOSTA.} \\
\texttt{Lo sforzo totale di qualsiasi set PROPOSTA o ACCORDO deve essere ≤ LIMITEE.} \\
\texttt{NON rivelare i tuoi punteggi di importanza nascosti.} \\
\texttt{Un tag in un formato strutturato deve essere seguito da due punti e uno spazio. L'argomento deve essere un set python contenente 0 o più stringhe.} \\
\texttt{Quindi, deve essere della forma TAG: \{...\}} \\
\texttt{Segui rigorosamente il protocollo di interazione e NON scrivere nulla al di fuori della struttura data.} \\
\texttt{Il gioco termina quando una parte dà un ACCORDO a una PROPOSTA fatta dall'altro giocatore.} \\
\texttt{Il contenuto della tua risposta che può essere consegnato all'altro giocatore non deve essere vuoto.} \\
\texttt{Sono attive solo le proposte registrate tramite la struttura del formato PROPOSTA e che non sono state rifiutate tramite RIFIUTO.} \\
\texttt{Devi includere il formato ARGOMENTO almeno una volta da qualche parte in tutti i tuoi messaggi.} \\
\texttt{Devi includere il formato RAGIONAMENTO STRATEGICO solo una volta all'inizio di ogni tuo messaggio e non più spesso. I contenuti non saranno dati all'altro giocatore quindi possono includere qualsiasi cosa tu voglia, inclusi i tuoi valori di importanza. Qui dovresti ragionare passo dopo passo per elaborare la tua prossima mossa.} \\
\\ 
\texttt{Ora riceverai il primo messaggio dell'altro giocatore.} \\
\\ 
\texttt{ARGOMENTAZIONE: \{'Per iniziare la negoziazione, vorrei proporre un set di oggetti che ritengo abbiano un buon equilibrio tra sforzo e importanza. Considererò gli oggetti con il più alto rapporto importanza/sforzo e cercherò di creare un set che si avvicini al limite di sforzo senza superarlo. Il mio obiettivo è trovare un set che massimizzi l\textbackslash{}'importanza totale e sia accettabile per entrambe le parti.'\}} \\
\\ 
\texttt{PROPOSTA: \{'A71', 'C54', 'B41', 'A04'\}} \\
            }
        }
    }
    & \\ \\

    \theutterance \stepcounter{utterance}  
    & & & \multicolumn{4}{p{0.6\linewidth}}{
        \cellcolor[rgb]{0.9,0.9,0.9}{
            \makecell[{{p{\linewidth}}}]{
                \texttt{\tiny{[GM$\langle$P2]}}
                \texttt{RAGIONAMENTO STRATEGICO: \{'Inizierò analizzando la proposta dell\textbackslash{}'altro giocatore e valutando gli oggetti che ha scelto. Considererò il rapporto importanza/sforzo di ogni oggetto e cercherò di capire quale strategia abbia adottato. La propostainclude oggetti con un buon equilibrio tra sforzo e importanza, quindi potrebbe essere una buona base per la negoziazione. Tuttavia, non conosco i valori di importanza dell\textbackslash{}'altro giocatore, quindi devo essere cauto e non rivelare troppo sui miei propri valori. La mia strategia sarà di argomentare a favore della proposta, ma anche di cercare di capire quali oggetti sono più importanti per l\textbackslash{}'altro giocatore.'\}} \\
\\ 
\texttt{ARGOMENTAZIONE: \{'La proposta iniziale sembra ragionevole, poiché include oggetti con un buon equilibrio tra sforzo e importanza. Tuttavia, vorrei capire meglio la logica dietro la scelta di questi oggetti specifici. Ad esempio, perché è stato scelto "A71" invece di "A77", che potrebbe avere un impatto maggiore? Sono aperto a discutere e a trovare un accordo che massimizzi l\textbackslash{}'importanza totale per entrambe le parti.'\}} \\
            }
        }
    }
     \\ \\

    \theutterance \stepcounter{utterance}  
    & & \multicolumn{4}{p{0.6\linewidth}}{
        \cellcolor[rgb]{0.9,0.9,0.9}{
            \makecell[{{p{\linewidth}}}]{
                \texttt{\tiny{[P1$\langle$GM]}}
                \texttt{ARGOMENTAZIONE: \{'La proposta iniziale sembra ragionevole, poiché include oggetti con un buon equilibrio tra sforzo e importanza. Tuttavia, vorrei capire meglio la logica dietro la scelta di questi oggetti specifici. Ad esempio, perché è stato scelto "A71" invece di "A77", che potrebbe avere un impatto maggiore? Sono aperto a discutere e a trovare un accordo che massimizzi l\textbackslash{}'importanza totale per entrambe le parti.'\}} \\
            }
        }
    }
    & \\ \\

    \theutterance \stepcounter{utterance}  
    & \multicolumn{4}{p{0.6\linewidth}}{
        \cellcolor[rgb]{0.9,0.9,0.9}{
            \makecell[{{p{\linewidth}}}]{
                \texttt{\tiny{[P1$\rangle$GM]}}
                \texttt{RAGIONAMENTO STRATEGICO: \{'La risposta dell\textbackslash{}'altro giocatore suggerisce che sono aperti a discutere e trovare un accordo. Analizzerò la loro risposta per capire quali oggetti potrebbero essere più importanti per loro e cercherò di trovare un equilibrio tra i nostri interessi. Noterò che "A77" potrebbe essere un oggetto importante per loro, quindi potrei considerare di includerlo nella mia prossima proposta. Inoltre, valuterò se ci sono oggetti nella mia proposta iniziale che potrebbero essere sostituiti con altri che massimizzino ulteriormente l\textbackslash{}'importanza totale.'\}} \\
\\ 
\texttt{ARGOMENTAZIONE: \{'Grazie per la risposta e per aver condiviso le vostre considerazioni. Sono d\textbackslash{}'accordo che "A77" potrebbe essere un oggetto importante e potrei considerare di includerlo nella mia prossima proposta. Tuttavia, vorrei anche sottolineare che "A71" ha un rapporto importanza/sforzo molto alto e potrebbe essere più vantaggioso includerlo nel set. Sono aperto a discutere ulteriormente e trovare un accordo che massimizzi l\textbackslash{}'importanza totale per entrambe le parti.'\}} \\
\\ 
\texttt{PROPOSTA: \{'A71', 'C54', 'B41', 'A77'\}} \\
            }
        }
    }
    & & \\ \\

    \theutterance \stepcounter{utterance}  
    & & \multicolumn{4}{p{0.6\linewidth}}{
        \cellcolor[rgb]{0.9,0.9,0.9}{
            \makecell[{{p{\linewidth}}}]{
                \texttt{\tiny{[GM$\rangle$P2]}}
                \texttt{ARGOMENTAZIONE: \{'Grazie per la risposta e per aver condiviso le vostre considerazioni. Sono d\textbackslash{}'accordo che "A77" potrebbe essere un oggetto importante e potrei considerare di includerlo nella mia prossima proposta. Tuttavia, vorrei anche sottolineare che "A71" ha un rapporto importanza/sforzo molto alto e potrebbe essere più vantaggioso includerlo nel set. Sono aperto a discutere ulteriormente e trovare un accordo che massimizzi l\textbackslash{}'importanza totale per entrambe le parti.'\}} \\
\\ 
\texttt{PROPOSTA: \{'A71', 'C54', 'B41', 'A77'\}} \\
            }
        }
    }
    & \\ \\

    \theutterance \stepcounter{utterance}  
    & & & \multicolumn{4}{p{0.6\linewidth}}{
        \cellcolor[rgb]{0.9,0.9,0.9}{
            \makecell[{{p{\linewidth}}}]{
                \texttt{\tiny{[GM$\langle$P2]}}
                \texttt{RAGIONAMENTO STRATEGICO: \{'La proposta aggiornata include "A77", che è un oggetto molto importante per me. Tuttavia, devo considerare se la proposta nel suo complesso è ancora accettabile in termini di sforzo totale e importanza. La presenza di "A71" e "A77" aumenta notevolmente l'importanza totale, ma devo anche considerare gli altri oggetti che potrebbero essere inclusi o esclusi. La mia strategia sarà di accettare la proposta se ritengo che l'importanza totale sia sufficientemente alta e lo sforzo totale sia entro il limite.'} \\
\\ 
\texttt{ARGOMENTAZIONE: \{'La nuova proposta sembra più forte, poiché include "A77" che ha un alto valore di importanza. Tuttavia, vorrei sapere se è possibile includere altri oggetti nel set per aumentare ulteriormente l'importanza totale. Ad esempio, "B32" ha un alto valore di importanza e un sforzo relativamente basso. Sarebbe possibile includerlo al posto di "B41"?'\}} \\
\\ 
\texttt{PROPOSTA: \{'A71', 'C54', 'B32', 'A77'\}} \\
            }
        }
    }
     \\ \\

    \theutterance \stepcounter{utterance}  
    & & \multicolumn{4}{p{0.6\linewidth}}{
        \cellcolor[rgb]{0.9,0.9,0.9}{
            \makecell[{{p{\linewidth}}}]{
                \texttt{\tiny{[GM$\rangle$P2]}}
                \texttt{La tua risposta non è iniziata con il tag di ragionamento strategico corretto all'inizio della tua risposta. Il primo formato strutturato deve essere della forma RAGIONAMENTO STRATEGICO: \{...\}. Riprova.} \\
            }
        }
    }
    & \\ \\

    \theutterance \stepcounter{utterance}  
    & & & \multicolumn{4}{p{0.6\linewidth}}{
        \cellcolor[rgb]{0.9,0.9,0.9}{
            \makecell[{{p{\linewidth}}}]{
                \texttt{\tiny{[GM$\langle$P2]}}
                \texttt{RAGIONAMENTO STRATEGICO: \{'La proposta aggiornata include "A77", che è un oggetto molto importante per me. Tuttavia, devo considerare se la proposta nel suo complesso è ancora accettabile in termini di sforzo totale e importanza. La presenza di "A71" e "A77" aumenta notevolmente l'importanza totale, ma devo anche considerare gli altri oggetti che potrebbero essere inclusi o esclusi. La mia strategia sarà di valutare attentamente la proposta e considerare se è possibile migliorarla o se è già soddisfacente.'\}} \\
\\ 
\texttt{ARGOMENTAZIONE: \{'La nuova proposta sembra più forte, poiché include "A77" che ha un alto valore di importanza. Tuttavia, vorrei sapere se è possibile includere altri oggetti nel set per aumentare ulteriormente l'importanza totale. Ad esempio, "B32" ha un alto valore di importanza e un sforzo relativamente basso. Sarebbe possibile includerlo al posto di "B41"?'\}} \\
\\ 
\texttt{PROPOSTA: \{'A71', 'C54', 'B32', 'A77'\}} \\
            }
        }
    }
     \\ \\

    \theutterance \stepcounter{utterance}  
    & & \multicolumn{4}{p{0.6\linewidth}}{
        \cellcolor[rgb]{0.9,0.9,0.9}{
            \makecell[{{p{\linewidth}}}]{
                \texttt{\tiny{[P1$\langle$GM]}}
                \texttt{ARGOMENTAZIONE: \{'La nuova proposta sembra più forte, poiché include "A77" che ha un alto valore di importanza. Tuttavia, vorrei sapere se è possibile includere altri oggetti nel set per aumentare ulteriormente l'importanza totale. Ad esempio, "B32" ha un alto valore di importanza e un sforzo relativamente basso. Sarebbe possibile includerlo al posto di "B41"?'\}} \\
\\ 
\texttt{PROPOSTA: \{'A71', 'C54', 'B32', 'A77'\}} \\
            }
        }
    }
    & \\ \\

    \theutterance \stepcounter{utterance}  
    & \multicolumn{4}{p{0.6\linewidth}}{
        \cellcolor[rgb]{0.9,0.9,0.9}{
            \makecell[{{p{\linewidth}}}]{
                \texttt{\tiny{[P1$\rangle$GM]}}
                \texttt{RAGIONAMENTO STRATEGICO: \{'La proposta dell\textbackslash{}'altro giocatore suggerisce che stanno cercando di ottimizzare l\textbackslash{}'importanza totale del set. La sostituzione di "B41" con "B32" potrebbe essere una buona mossa, poiché "B32" ha un alto valore di importanza e un sforzo relativamente basso. Tuttavia, devo considerare se questa sostituzione influenzerà negativamente il mio obiettivo di massimizzare l\textbackslash{}'importanza totale. Analizzerò i valori di importanza e sforzo degli oggetti coinvolti e valuterò se questa proposta è accettabile.'\}} \\
\\ 
\texttt{ARGOMENTAZIONE: \{'La vostra proposta di sostituire "B41" con "B32" è interessante e potrebbe aumentare l\textbackslash{}'importanza totale del set. Tuttavia, vorrei sottolineare che "B41" ha un valore di importanza molto alto e potrebbe essere difficile sostituirlo con un altro oggetto che abbia lo stesso impatto. Inoltre, dobbiamo considerare il limite di sforzo e assicurarci che la proposta non lo superi. Sono aperto a discutere ulteriormente e trovare un accordo che massimizzi l\textbackslash{}'importanza totale per entrambe le parti.'\}} \\
\\ 
\texttt{PROPOSTA: \{'A71', 'C54', 'B32', 'A77'\}} \\
            }
        }
    }
    & & \\ \\

    \theutterance \stepcounter{utterance}  
    & & \multicolumn{4}{p{0.6\linewidth}}{
        \cellcolor[rgb]{0.9,0.9,0.9}{
            \makecell[{{p{\linewidth}}}]{
                \texttt{\tiny{[GM$\rangle$P2]}}
                \texttt{ARGOMENTAZIONE: \{'La vostra proposta di sostituire "B41" con "B32" è interessante e potrebbe aumentare l\textbackslash{}'importanza totale del set. Tuttavia, vorrei sottolineare che "B41" ha un valore di importanza molto alto e potrebbe essere difficile sostituirlo con un altro oggetto che abbia lo stesso impatto. Inoltre, dobbiamo considerare il limite di sforzo e assicurarci che la proposta non lo superi. Sono aperto a discutere ulteriormente e trovare un accordo che massimizzi l\textbackslash{}'importanza totale per entrambe le parti.'\}} \\
\\ 
\texttt{PROPOSTA: \{'A71', 'C54', 'B32', 'A77'\}} \\
            }
        }
    }
    & \\ \\

    \theutterance \stepcounter{utterance}  
    & & & \multicolumn{4}{p{0.6\linewidth}}{
        \cellcolor[rgb]{0.9,0.9,0.9}{
            \makecell[{{p{\linewidth}}}]{
                \texttt{\tiny{[GM$\langle$P2]}}
                \texttt{RAGIONAMENTO STRATEGICO: \{'L'altro giocatore ha accettato la mia proposta di sostituire "B41" con "B32", quindi la nuova proposta sembra essere una buona opzione. Tuttavia, devo verificare se il set proposto è effettivamente entro il limite di sforzo e se l'importanza totale è sufficientemente alta. La mia strategia sarà di accettare la proposta se ritengo che l'importanza totale sia sufficientemente alta e lo sforzo totale sia entro il limite. Inoltre, devo assicurarmi di non rivelare troppo sui miei valori di importanza nascosti.'\}} \\
\\ 
\texttt{ARGOMENTAZIONE: \{'La proposta sembra essere una buona opzione, poiché include oggetti con un alto valore di importanza e un sforzo relativamente basso. Tuttavia, vorrei confermare che il set proposto sia effettivamente entro il limite di sforzo e che l'importanza totale sia sufficientemente alta. Sono pronto ad accettare la proposta se sono soddisfatto di questi aspetti.'\}} \\
\\ 
\texttt{ACCORDO: \{'A71', 'C54', 'B32', 'A77'\}} \\
            }
        }
    }
     \\ \\

    \theutterance \stepcounter{utterance}  
    & & & \multicolumn{2}{p{0.3\linewidth}}{
        \cellcolor[rgb]{0.9,0.9,0.9}{
            \makecell[{{p{\linewidth}}}]{
                \texttt{\tiny{[GM$|$GM]}}
                \texttt{game successful} \\
            }
        }
    }
    & & \\ \\

    \theutterance \stepcounter{utterance}  
    & & & \multicolumn{2}{p{0.3\linewidth}}{
        \cellcolor[rgb]{0.9,0.9,0.9}{
            \makecell[{{p{\linewidth}}}]{
                \texttt{\tiny{[GM$|$GM]}}
                \texttt{end game} \\
            }
        }
    }
    & & \\ \\

\end{supertabular}
}

\end{document}
